\documentclass[10pt]{article}
\usepackage[english]{babel}
\usepackage{../../../meta-inf/lib/naproche}
\usepackage{amssymb}
\usepackage{mathtools} % for \coloneq

\usepackage{stex-highlighting}
\providebool{emph} % "\newbool{emph}" does not work...
\setbool{emph}{false}
\colorlet{emphcolor}{violet}
\let\oldemph\emph
\renewcommand\emph[1]{\setbool{emph}{true}\ifbool{forthel}{\textcolor{emphcolor}{\itshape#1}}{\oldemph{#1}}\setbool{emph}{false}}
\renewcommand{\varemph}[1]{\ifbool{emph}{\textcolor{emphcolor}{#1}}{\textcolor{black}{#1}}}

\usepackage[right=6cm,left=3cm,bottom=3cm,marginparwidth=5cm]{geometry}

\usepackage{fancyhdr}
\renewcommand{\sectionmark}[1]{\markboth{#1}{}} 
\def\libarchive{}
\pagestyle{fancy}
\fancyhead[L]{\libarchive}
\fancyhead[C]{\nouppercase\leftmark}  % section title
\fancyhead[R]{\thepage}               % page number
\fancyfoot[C]{}                       % No page number in footer

\usepackage[nobottomtitles]{titlesec}
\titlespacing*{\section}{0pt}{30pt}{0pt}
\titlespacing*{\subsection}{0pt}{30pt}{0pt}
\titlespacing*{\subsubsection}{0pt}{30pt}{0pt}

\documentclass[12pt,oneside]{book}

\usepackage[foundations]{../../lib/tex/naproche}
\usepackage{../../lib/tex/libraries}
\usepackage{graphicx}
\usepackage{float}
\usepackage{caption}
\usepackage{footnote}

\makesavenoteenv{tabular} % Make footnotes work in tabular environments


\title{Foundations of Mathematics}
\author{Marcel Schütz}
\date{2022}

\begin{document}
  \maketitle

  \tableofcontents

  \begin{figure}[H]
    \centering
    \fbox{\includegraphics[width=0.9\linewidth]{./dependency-graph/graph.png}}
    \caption*{Interdependencies of the chapters}
  \end{figure}


  \section*{Introduction}

  This is a library providing a foundation of mathematics based on a
  Kelley-Morse like class theory with urelements.
  It introduces common operations on classes like unions or intersections
  (\cref{chapter:classes}) together with detailed proofs of their algebraic
  properties (\cref{chapter:computation-laws-for-classes}), the symmetric
  difference of two classes (\cref{chapter:symmetric-difference}) and the
  notions of ordered pairs and Cartesian products
  (\cref{chapter:pairs-and-products}) as well as proofs of the algebraic
  properties of the latter (\cref{chapter:computation-laws-for-products}).
  Moreover, it provides common operations on maps (\cref{chapter:maps}), various
  properties of images and preimages (\cref{chapter:image-and-preimage}) and the
  notions of injectivity, surjectivity, bijectivity
  (\cref{chapter:injections-surjections-bijections}) and invertibility of maps
  (\cref{chapter:invertible-maps}).
  The library provides an axiom system characterizing sets (\cref{chapter:sets})
  and, furthermore, it covers the notions of binary relations
  (\cref{chapter:binary-relations}), fixed-points of subset preserving maps
  (\cref{chapter:fixed-points}), including and equinumerosity
  (\cref{chapter:equinumerosity}).

  As two famous results it includes the Knaster-Tarski fixed point theorem
  (\cref{FOUNDATIONS_12_8420450166112256}) and the Cantor-Schröder-Bernstein
  theorem (\cref{FOUNDATIONS_13_1913663275401216}).

  \paragraph*{Usage.}
  At the very beginning of each chapter you can find the name of its source
  file, e.g. \path{foundations/sections/01_classes.ftl.tex} for
  \cref{chapter:classes}. This filename can be used to import the chapter via
  \Naproche's \texttt{readtex} instruction to another ForTheL text, e.g.:
  \begin{center}
    \verb`[readtex \path{foundations/sections/01_classes.ftl.tex}]`
  \end{center}

  \paragraph*{Checking times.}
  The checking times for each of the chapters may vary from computer to
  computer, but on mid-range hardware they are likely to be similar to those
  given in table below:

  \begin{center}
    \begin{tabular}{c|c|c}

      & \multicolumn{2}{c}{\textbf{Checking time}}
      \\
      \textbf{Chapter}
      & \textbf{without dependencies}     & \textbf{with dependencies}
      \\ \hline
      \ref{chapter:classes}
      & 00:04 min                         & 00:04 min
      \\
      \ref{chapter:computation-laws-for-classes}
      & 00:12 min                         & 00:16 min
      \\
      \ref{chapter:symmetric-difference}
      & 00:32 min                         & 00:48 min
      \\
      \ref{chapter:pairs-and-products}
      & 00:08 min                         & 00:12 min
      \\
      \ref{chapter:computation-laws-for-products}
      & 01:36 min                         & 01:56 min
      \\
      \ref{chapter:maps}
      & 01:13 min                         & 01:25 min
      \\
      \ref{chapter:image-and-preimage}
      & 01:28 min                         & 02:53 min
      \\
      \ref{chapter:injections-surjections-bijections}
      & 00:38 min                         & 02:03 min
      \\
      \ref{chapter:invertible-maps}
      & 02:20 min                         & 04:23 min
      \\
      \ref{chapter:sets}
      & 02:17 min                         & 06:40 min
      \\
      \ref{chapter:binary-relations}
      & 00:14 min                         & 06:54 min
      \\
      \ref{chapter:fixed-points}
      & 00:33 min                         & 07:13 min
      \\
      \ref{chapter:equinumerosity}
      & 01:48 min                         & 09:01 min
    \end{tabular}
  \end{center}


  \subfile{sections/01_classes.ftl.tex}
  \subfile{sections/02_computation-laws-for-classes.ftl.tex}
  \subfile{sections/03_symmetric-difference.ftl.tex}
  \subfile{sections/04_pairs-and-products.ftl.tex}
  \subfile{sections/05_computation-laws-for-products.ftl.tex}
  \subfile{sections/06_maps.ftl.tex}
  \subfile{sections/07_image-and-preimage.ftl.tex}
  \subfile{sections/08_injections-surjections-bijections.ftl.tex}
  \subfile{sections/09_invertible-maps.ftl.tex}
  \subfile{sections/10_sets.ftl.tex}
  \subfile{sections/11_binary-relations.ftl.tex}
  \subfile{sections/12_fixed-points.ftl.tex}
  \subfile{sections/13_equinumerosity.ftl.tex}
\end{document}

\usepackage{amssymb}

\newcommand{\Nat}{\mathbb{N}}
\newcommand{\Prime}{\mathbb{P}}
\renewcommand{\succ}{\textrm{succ}}
\newcommand{\pred}{\textrm{pred}}
\newcommand{\add}{\textrm{add}}
\newcommand{\mul}{\textrm{mul}}
\renewcommand{\exp}{\textrm{exp}}
\newcommand{\fac}{\textrm{fac}}
\renewcommand{\div}{\mathop{\textrm{div}}}
\renewcommand{\mod}{\mathop{\textrm{mod}}}

\begin{document}
  \begin{imports}
    \begin{forthel}
      %[prove off][check off]
      [readtex \path{libraries/source/arithmetics/multiplication-and-ordering.ftl.tex}]
      %[prove on][check on]
    \end{forthel}
  \end{imports}


  \section*{Modular Arithmetics}

  \subsection*{Quotients and Remainders}

  \begin{forthel}
    \begin{definition}[id=ARITHMETIC_08_8621463798022144,printid]
      Let $n, m$ be natural numbers such that $m \neq 0$.
      $n \div m$ is the natural number $q$ such that $n = (m \cdot q) + r$ for some natural number $r$ that is less than $m$.
    \end{definition}

    Let the quotient of $n$ over $m$ stand for $n \div m$.
  \end{forthel}

  \begin{forthel}
    \begin{definition}[id=ARITHMETIC_08_3560980160184320,printid]
      Let $n, m$ be natural numbers such that $m \neq 0$.
      $n \mod m$ is the natural number $r$ such that $r < m$ and there exists a natural number $q$ such that $n = (m \cdot q) + r$.
    \end{definition}

    Let the remainder of $n$ over $m$ stand for $n \mod m$.
  \end{forthel}


  \subsection*{Basic Properties}

  \begin{forthel}
    \begin{definition}[id=ARITHMETIC_08_5448561831444480,printid]
      Let $n, m, k$ be natural numbers such that $k \neq 0$.
      $n \equiv m \pmod{k}$ iff $n \mod k = m \mod k$.
    \end{definition}

    Let $n$ and $m$ are congruent modulo $k$ stand for $n \equiv m \pmod{k}$.
  \end{forthel}

  \begin{forthel}
    \begin{proposition}[id=ARITHMETIC_08_3818318544764928,printid]
      Let $n, k$ be natural numbers such that $k \neq 0$.
      Then $n \equiv n \pmod{k}$.
    \end{proposition}
    \begin{proof}
      We have $n \mod k = n \mod k$.
      Hence $n \equiv n \pmod{k}$.
    \end{proof}
  \end{forthel}

  \begin{forthel}
    \begin{proposition}[id=ARITHMETIC_08_2337210737098752,printid]
      Let $n, m, k$ be natural numbers such that $k \neq 0$.
      If $n \equiv m \pmod{k}$ then $m \equiv n \pmod{k}$.
    \end{proposition}
    \begin{proof}
      Assume $n \equiv m \pmod{k}$.
      Then $n \mod k = m \mod k$.
      Hence $m \mod k = n \mod k$.
      Thus $m \equiv n \pmod{k}$.
    \end{proof}
  \end{forthel}

  \begin{forthel}
    \begin{proposition}[id=ARITHMETIC_08_7464329746055168,printid]
      Let $n, m, l, k$ be natural numbers such that $k \neq 0$.
      If $n \equiv m \pmod{k}$ and $m \equiv l \pmod{k}$ then $n \equiv l \pmod{k}$.
    \end{proposition}
    \begin{proof}
      Assume $n \equiv m \pmod{k}$ and $m \equiv l \pmod{k}$.
      Then $n \mod k = m \mod k$ and $m \mod k = l \mod k$.
      Hence $n \mod k = l \mod k$.
      Thus $n \equiv l \pmod{k}$.
    \end{proof}
  \end{forthel}

  \begin{forthel}
    \begin{proposition}[id=ARITHMETIC_08_2034122983735296,printid]
      Let $n, m, k$ be natural numbers such that $k \neq 0$.
      Assume $n \geq m$.
      Then $n \equiv m \pmod{k}$ iff $n = (k \cdot x) + m$ for some natural number $x$.
    \end{proposition}
    \begin{proof}
      Case $n \equiv m \pmod{k}$.
        Then $n \mod k = m \mod k$.
        Take a natural number $r$ such that $r < k$ and $n \mod k = r = m \mod k$.
        Take a nonzero natural number $l$ such that $k = r + l$.
        Consider natural numbers $q,q'$ such that $n = (q \cdot k) + r$ and $m = (q' \cdot k) + r$.

        Then $q \geq q'$. \\
        Proof.
          Assume the contrary.
          Then $q < q'$.
          Hence $q \cdot k < q' \cdot k$.
          Thus $(q \cdot k) + r < (q' \cdot k) + r$ (by \printref{ARITHMETIC_04_7354062662008832}).
          Indeed $q \cdot k$ and $q' \cdot k$ are natural numbers.
          Therefore $n < m$.
          Contradiction.
        Qed.

        Take a natural number $x$ such that $q = q' + x$.

        Let us show that $n = (k \cdot x) + m$.
          We have
          \[  (k \cdot x) + m                       \]
          \[    = (k \cdot x) + ((q' \cdot k) + r)  \]
          \[    = ((k \cdot x) + (q' \cdot k)) + r  \]
          \[    = ((k \cdot x) + (k \cdot q')) + r  \]
          \[    = (k \cdot (q' + x)) + r            \]
          \[    = (k \cdot q) + r                   \]
          \[    = n.                                \]
        End.
      End.

      Case $n = (k \cdot x) + m$ for some natural number $x$.
        Consider a natural number $x$ such that $n = (k \cdot x) + m$.
        Take natural numbers $r, r'$ such that $n \mod k = r$ and $m \mod k = r'$.
        Then $r, r' < k$.
        Take natural numbers $q, q'$ such that $n = (k \cdot q) + r$ and $m = (k \cdot q') + r'$.
        Then
        \[  (k \cdot q) + r                         \]
        \[    = n                                   \]
        \[    = (k \cdot x) + m                     \]
        \[    = (k \cdot x) + ((k \cdot q') + r')   \]
        \[    = ((k \cdot x) + (k \cdot q')) + r'   \]
        \[    = (k \cdot (x + q')) + r'.            \]
        Hence $r = r'$.
        Thus $n \mod k = m \mod k$.
        Therefore $n \equiv m \pmod{k}$.
      End.
    \end{proof}
  \end{forthel}

  \begin{forthel}
    \begin{proposition}[id=ARITHMETIC_08_2988318228742144,printid]
      Let $n, m, k, k'$ be natural numbers such that $k, k' \neq 0$.
      If $n \equiv m \pmod{k \cdot k'}$ then $n \equiv m \pmod{k}$.
    \end{proposition}
    \begin{proof}
      Assume $n \equiv m \pmod{k \cdot k'}$.

      Case $n \geq m$.
        We can take a natural number $x$ such that $n = ((k \cdot k') \cdot x) + m$.
        Then $n = (k \cdot (k' \cdot x)) + m$.
        Hence $n \equiv m \pmod{k}$.
      End.

      Case $m \geq n$.
        We have $m \equiv n \pmod{k \cdot k'}$.
        Hence we can take a natural number $x$ such that $m = ((k \cdot k') \cdot x) + n$.
        Then $m = (k \cdot (k' \cdot x)) + n$.
        Thus $m \equiv n \pmod{k}$.
        Therefore $n \equiv m \pmod{k}$.
      End.
    \end{proof}
  \end{forthel}

  \begin{forthel}
    \begin{corollary}[id=ARITHMETIC_08_5895145169879040,printid]
      Let $n, m, k, k'$ be natural numbers such that $k, k' \neq 0$.
      If $n \equiv m \pmod{k \cdot k'}$ then $n \equiv m \pmod{k'}$.
    \end{corollary}
    \begin{proof}
      Assume $n \equiv m \pmod{k \cdot k'}$.
      Then $n \equiv m \pmod{k' \cdot k}$.
      Hence $n \equiv m \pmod{k'}$.
    \end{proof}
  \end{forthel}

  \begin{forthel}
    \begin{proposition}[id=ARITHMETIC_08_5984712287846400,printid]
      Let $n, k$ be natural numbers such that $k \neq 0$.
      Then $n + k \equiv n \pmod{k}$.
    \end{proposition}
    \begin{proof}
      Take $r = n \mod k$ and $r' = (n + k) \mod k$.
      Consider a $q \in \Nat$ such that $n = (k \cdot q) + r$ and $r < k$.
      Consider a $q' \in \Nat$ such that $n + k = (k \cdot q') + r'$ and
      $r' < k$.
      Then $(k \cdot q') + r'
        = n + k
        = ((k \cdot q) + r) + k
        = (k + (k \cdot q)) + r
        = (k \cdot (q + 1)) + r$.
      Hence $r = r'$.
      Consequently $n \mod k = (n + k) \mod k$.
      Thus $n + k \equiv n \pmod{k}$.
    \end{proof}
  \end{forthel}
\end{document}
